\section{Ánh xạ}

Ánh xạ là một tên gọi khác của hàm. Giống như định nghĩa của hàm trong giải tích vị từ bậc nhất, ánh xạ cũng có tính chất cơ bản là ứng với mỗi hạng thức hoặc nhóm hạng thức, ánh xạ sẽ cho ra và chỉ cho ra một đối tượng. Cách hiểu của ánh xạ này là tương đối mới, bởi từ trước thế kỉ \Rmnum{19}, ánh xạ được xem là một đối tượng được xác định bởi một công thức nào đó. Định nghĩa này thì đơn giản hơn và có tính ứng dụng tốt hơn định nghĩa hiện nay, nhưng đi kèm với đó là những sự tranh cãi về giới hạn trong việc mô tả công thức. Ngoài ra, hai ánh xạ có thể có định nghĩa khác nhau nhưng lại có thể hoàn toàn đồng nhất. Định nghĩa cũ vừa có tính dư thừa vừa không đủ bao quát, dẫn đến nhu câu đưa ra định nghĩa mới cho hàm, hay ánh xạ.

\subsection{Định nghĩa ánh xạ}

Một \defText{ánh xạ} hay \defText{hàm số}\footnote{Cũng có thể gọi là \defText{hàm}, nhưng để tránh bị lặp với khái niệm hàm tử giải tích vị từ bậc nhất, tác giả sẽ không sử dụng tên ``hàm''.} $f$ từ $x$ đến $y$ là một quan hệ $f$ từ $x$ đến $y$ thỏa mãn \begin{equation}
   \defMath{\forall a \in x, \exists ! b \in y, a\;f\;b }.
   \label{eq:li_thuyet_tap_hop:tien_de:anh_xa:dinh_nghia}
\end{equation}
Kí hiệu thường xuyên sử dụng thay vì $a\;f\;b$ là $\defMath{b = f(a)}$. $b$ khi này được gọi là \defText{ảnh} của $a$ qua $f$ và $a$ được gọi là $\defText{tạo ảnh}$ của $b$ bởi $f$. Có một điểm cần lưu ý, đó là từ định nghĩa ánh xạ, chúng ta cần phải có tập xác định của $f$ thì phải bằng tập nguồn $x$. Ngoài ra, nếu chiếu theo định nghĩa này thì tập nguồn $x$ có thể rỗng, bởi khi đó $\forall a, \neg a \in x$ khiến cho \refeq{eq:li_thuyet_tap_hop:tien_de:anh_xa:dinh_nghia} luôn đúng. Tác giả sẽ chấp nhận tập nguồn rỗng, nhưng không phải tài liệu nào cũng sẽ đồng tình với tác giả.

Từ $x$ và $y$, xây dựng tích đề các $x \times y$. Do mọi quan hệ đơn giản chỉ là tập con của $x \times y$ nên $\mathfrak{P}(x \times y)$ sẽ chứa toàn bộ các quan hệ trên $x \times y$. Và bởi vì mọi ánh xạ là quan hệ, chúng ta có tập hợp các ánh xạ từ $x$ đến $y$ là
$$
   \defMath{y^x = \big\{z \in \mathfrak{P}(x \times y)\;\big|\;z \text{\defText{~là một ánh xạ}} \big\}}.
$$

Vì ánh xạ là một khái niệm thường xuyên được sử dụng, người ta đã sáng tạo nhiều cách để kí hiệu nó mà chứa nhiều thông tin nhất có thể. Đây là một trong những cách kí hiệu thường dùng:
\[
\begin{array}{rcl}
\defMath{f: x} & \defMath{\to} & \defMath{y} \\
   \defMath{z} & \defMath{\mapsto} & \defMath{f(z)}
\end{array}.
\]

Lấy một số ví dụ về ánh xạ:
\begin{itemize}
   \item Với tập hợp $A$ bất kì, kí hiệu $\begin{array}[t]{rcl}
\defMath{\mathds{1}_A: A} & \defMath{\to} & \defMath{A} \\
   \defMath{z} & \defMath{\mapsto} & \defMath{z}
\end{array}$ là \defText{ánh xạ đồng nhất} của $A$. Ánh xạ này đôi khi còn được kí hiệu là $\defMath{\Id_A}$ hoặc $\defMath{\id_A}$.
   \item Cho $A$ là một tập hợp và $E \in \mathfrak{P}(A)$. \defText{Ánh xạ nhúng chính tắc} từ $E$ vào $A$ là ánh xạ xác định bởi:\begin{equation*}
      \begin{array}[t]{rcl}
      \defMath{\mathbf{i}_{E; A}: E} & \defMath{\to} & \defMath{A} \\
         \defMath{z} & \defMath{\mapsto} & \defMath{z}
      \end{array}.
   \end{equation*}
   \item Cho $A$ và $E$ là hai tập hợp và hạng $x \in E$. \defText{Ánh xạ hằng} $x$ là ánh xạ, thường được kí hiệu trùng với hạng thức, có định nghĩa $\begin{array}[t]{rcl}
      \defMath{x: A} & \defMath{\to} & \defMath{E} \\
         \defMath{z} & \defMath{\mapsto} & \defMath{x}
      \end{array}.$
\end{itemize}

Cho một ánh xạ $f: x \to y$, $A$ là một tập con của $x$. Gọi $g$ là quan hệ từ $A$ đến $y$ thỏa mãn $a\;g\;b \iff f(a) = b$. Với mỗi $a$, $b = f(a)$ tồn tại và là duy nhất theo định nghĩa ánh xạ, do đó $g$ cũng xác định một ánh xạ. Một cách cảm quan, nếu $a \in A$ thì $f(a) = g(a)$. Có thể thấy, $g$ là một \defText{ánh xạ cảm sinh} (hay \defText{ánh xạ thu hẹp}) từ $x$ xuống $A$, và ngược lại, $f$ là một \defText{ánh xạ thác triển} (hay \defText{ánh xạ mở rộng}) từ $A$ ra $x$.

Cuối cùng, xét $y$ là một tập con của $x$ và một ánh xạ $f: x \to x$. $y$ được gọi là \defText{ổn định} đối với $f$ khi và chỉ khi \begin{equation*}
   \defMath{\forall z \in y, f(z) \in y}.
\end{equation*}
Tính chất ổn định có thể được mở rộng ra thành các khái niệm đối xứng, thể hiện tập hợp các điểm trong $y$ vẫn giữ nguyên tính chất vật lí khi thực hiện phép biến đổi này.

\exercise Cho $A$ và $B$ là hai tập hợp sao cho tồn tại ánh xạ $f: A \to B$.
\begin{enumerate}
   \item Chứng minh rằng nếu $B = \varnothing$ thì $A = \varnothing$.
   \item Chứng minh rằng nếu $\left[\begin{aligned}
      A = \varnothing \\ B = \varnothing
   \end{aligned}\right.$ thì $f = \varnothing$.
\end{enumerate}

\solution

\resetSubexercise Giả sử $A \neq \varnothing$, nghĩa là tồn tại $x \in A$. Khi này, do tồn tại ánh xạ $f$, theo định nghĩa ánh xạ, tồn tại duy nhất một $y \in B$ sao cho $x\;f\;y$. Tồn tại duy nhất một phần tử để thỏa mãn điều kiện gì đó có nghĩa là trước hết phần tử đó phải tồn tại. Do đó, $B$ không rỗng.

Vì nếu $A \neq \varnothing$ thì $B \neq \varnothing$, nên, chúng ta cũng có mệnh đề tương đương
\begin{equation*}
   B = \varnothing \implies A = \varnothing.
\end{equation*}
Điều phải chứng minh.

\stepSubexercise $\begin{cases}
   B = \varnothing \implies A = \varnothing \\
   A = \varnothing \implies A = \varnothing
\end{cases}$ nên $\left[\begin{aligned}
   A = \varnothing \\ B = \varnothing
\end{aligned}\right. \implies A = \varnothing$. 

\exercise Cho $A$ và $B$ là hai tập hợp (và $B$ khác rỗng) sao cho tồn tại ánh xạ $f: A \to B$. Nếu $B \neq \varnothing$, gọi $b$ là một phần tử của $B$ và tập $C$ thỏa mãn $C \supseteq A$. Xây dựng quan hệ $g$ từ $C$ đến $B$ sao cho \begin{equation*}
   \forall x \in C, \forall y \in B, \left(x\;g\;y \iff \left[\begin{aligned}         
      x \in A &\land y = f(x) \\
      x \notin A &\land y = b
\end{aligned}
\right.\right).
\end{equation*} 
Chứng minh rằng $g$ là ánh xạ thác triển của $f$.

\solution

Xét một phần tử $c \in C$. Có hai trường hợp có thể xảy ra,
\begin{itemize}
   \item Trường hợp \coloringCrimson{$c \in A$}: Do $f: A \to B$ là một ánh xạ và $c \in A$, tồn tại một phần tử $y_c \in B$ sao cho $y_c = f(c)$. Theo định nghĩa của $g$, $c \in A$ và $y_c = f(c)$ cho $c\;g\;y_c$. Vậy, xác định được $y_c$ để $c\;g\;y_c$.
   
   Tiếp theo, chúng ta chứng minh tại sao giá trị $y_c$ này là duy nhất. Giả sử có thêm $y_{cc}$ thỏa mãn $c\;g\;y_{cc}$. Theo định nghĩa $g$, chúng ta có $
      \left[\begin{aligned}
         c \in A &\land y_{cc} = f(c) \\
         c \notin A &\land y_{cc} = b
      \end{aligned}\right..
   $ Do $c \in A$ nên $\overline{c \notin A \land y = b}$ và kéo theo $y_{cc} = f(c) = y_c$. Và vì thế, $y_c$ là duy nhất.

   \item Trường hợp \coloringDodgerblue{$c \notin A$}: Vì $c \notin A$ và $b = b$, nên $c\;g\;b$. 
   
   Để chứng minh $c$ chỉ có mối quan hệ duy nhất với $b$, tương tự như trước, giả sử có $\yNom$ nữa sao cho $c\;g\;\yNom$. Chúng ta có $\neg c \in A$ nên $\overline{c \in A \land \yNom = f(c)}$, dẫn đến phải có $\yNom = b$.
\end{itemize}

Trong cả hai trường hợp, với mọi $c \in C$, tồn tại duy nhất một phần tử $y \in B$ sao cho $c\;g\;y$. Theo định nghĩa ánh xạ, $g$ là một ánh xạ từ $C$ đến $B$.

Với mọi $x \in A$, $f(x) = g(x)$ theo chứng minh ở trên. Do đó, có thể viết $f$ dưới dạng:
\begin{equation*}
   \begin{array}[t]{rcl}
      f: A & \to & B \\
      x & \mapsto & g(x)
   \end{array}.
\end{equation*}
Có $A\subseteq C$, từ định nghĩa, chúng ta có $g$ là ánh xạ thác triển của $f$. Điều phải chứng minh.

\exercise Cho $A$ và $B$ là hai tập hợp. Xác định hai ánh xạ $f: A \to B$ và $g: A \to B$ sao cho 
\begin{equation*}
   \forall x \in A, f(x) = g(x). 
\end{equation*}
Chứng minh rằng $f = g$.

\solution

Với mọi $a$ là một phần tử thuộc $f$, tồn tại $b \in A$ và $c \in B$ sao cho $a = (b; c)$. Có 
\begin{equation*}
   c = f(b) = g(b)
\end{equation*} 
nên $(b; c)$ cũng thuộc $g$. Vì vậy, 
$$
   \forall a \in f, a \in g.
$$
Tương tự, $\forall a \in g, a \in f$. Do đó, $f = g$. Điều phải chứng minh.

\exercise Cho $A$ và $B$ là hai tập hợp. $C$ có định nghĩa là
\begin{equation*}
   C = \left\{x \in \mathfrak{P}(A) \times (A \times B)\;\middle|\;\exists y \in \mathfrak{P}(A), \exists z, \left(z \text{~là hàm số từ~}y\text{~sang~}B \land x = (y; z)\right)\right\}.
\end{equation*}
Định nghĩa trong $C$ quan hệ $R$ sao cho với mọi đôi có thứ tự $(c; d)$ và $(c'; d')$ trong $C$,
\begin{equation*}
   (c; d)\;R\;(c'; d') \iff \begin{cases}
      c \subseteq c' \\
      \forall e \in c, d(e) = d'(e) 
   \end{cases}.
\end{equation*}

\begin{enumerate}
   \item Chứng minh rằng $R$ là một quan hệ thứ tự trên $C$.
   \item Xác định\footnote{Với những bài tập yêu cầu \emph{xác định} hay \emph{tìm} giá trị của một hoặc nhiều đối tượng (hay còn được gọi tắt là \emph{tìm nghiệm}), thông thường, chúng ta cần phải liệt kê các giá trị thỏa mãn hoặc chỉ ra phương pháp để có thể \emph{dễ dàng} xây dựng nên các nghiệm đáp ứng điều kiện đã cho. Không có một định nghĩa khách quan để giải nghĩa thế nào là ``dễ dàng''. Chẳng hạn, với câu hỏi tìm $x$ để $x \cup A = B$, chúng ta có thể trả lời là:
\begin{center}
   ``Tập $x$ thỏa mãn là những tập có tính chất $x \cup A = B$.''
\end{center} nhưng trả lời như vậy thì rõ ràng là không thỏa đáng. Một câu trả lời phải có giá trị toán học hoặc ứng dụng ở một mức độ nào đó. Theo quan niệm của tác giả, khi giải một bài toán tìm nghiệm, cần phải chỉ ra một thủ tục tính toán, sử dụng các giá trị cho trước và phép tính đã định nghĩa để xây dựng được các giá trị thỏa mãn. Trong đó, liệt kê là một thủ tục có thể sử dụng khi có ít nghiệm hợp lệ.} các phần tử cực đại và cực tiểu của $C$ (theo quan hệ $R$).
\end{enumerate}

\solution

\resetSubexercise Khẳng định 
\begin{equation*}
   \begin{cases}
      c \subseteq c \\
      \forall e \in c, d(e) = d(e)
   \end{cases}
\end{equation*}
là hiển nhiên. Do đó, $(c; d)\; R \; (c; d)$ luôn đúng với mọi $f = (c; d)$ trong $C$, hay $\forall f \in C, f\;R\;f$. $R$ là một quan hệ có tính phản xạ.

Xét $g$ và $h$ thuộc $C$ sao cho $g\;R\;h$ và $h\;R\;g$. Tồn tại $a_g$, $a_h$ thuộc $A$, $b_g$, $b_h$ thuộc $B$ sao cho $g = (a_g; b_g)$ và $h = (a_h; b_h)$. Có
$$\left(g\;R\;h \iff \begin{cases}
   a_g \subseteq a_h \\
   \forall e \in a_g, b_g(e) = b_h(e)
\end{cases}\right) \land 
\left(h\;R\;g \iff \begin{cases}
   a_h \subseteq a_g \\
   \forall e \in a_h, b_h(e) = b_g(e)
\end{cases} \right)$$
nên suy ra $a_g = a_h$. Lấy $p_g = (i_g; o_g)$ với $p_g \in b_g$ bất kì. Do $\forall e \in a_g, b_g(e) = b_h(e)$, \begin{equation*}
   b_g(i_g) = b_h(i_g),
\end{equation*} hay $o_g = b_h(i_g)$. Điều này tương đương với $p_g \in b_h$, và vì thế $b_g \subseteq b_h$. Quan hệ bao hàm ngược lại $b_g \supseteq b_h$ cũng hiển nhiên. Do đó $b_g = b_h$. Vì thế, $(a_g; b_g) = (a_h; b_h)$ hay $g = h$. Chúng ta thấy $R$ phản đối xứng.

Xét $\xNom$, $\yNom$ và $\zNom$ trong $C$ để $\xNom\;R\;\yNom$ và $\yNom\;R\;\zNom$. Theo định nghĩa của $R$, chúng ta có $\xNom = (m; n)$, $\yNom = (p; q)$, và $\zNom = (s; t)$ thỏa mãn
\begin{center}
   $\begin{cases}
      m \subseteq p \\
      \forall e \in m, n(e) = q(e)
   \end{cases}$ và $\begin{cases}
      p \subseteq s \\
      \forall e \in p, q(e) = t(e)
   \end{cases}$.
\end{center}
Theo tính chất bắc cầu của $\subseteq$, $m \subseteq s$. Ngoài ra, $m \subseteq p \implies \forall e \in m, e \in p$. Kết hợp với $\forall e \in p, q(e) = t(e)$, do $\implies$ cũng có tính bắc cầu nên
\begin{equation*}
   \forall e \in m, q(e) = t(e).
\end{equation*}
Đã có $\forall e \in m, n(e) = q(e)$, dễ thấy rằng \begin{equation*}
   \forall e \in m, n(e) = t(e).
\end{equation*}

Vì $m \subseteq s$ và $\forall e \in m, n(e) = t(e)$, theo định nghĩa $R$, $\xNom\;R\;\zNom$. Do $\xNom$, $\yNom$ và $\zNom$ được lấy bất kì, chúng ta khẳng định được rằng $R$ có tính chất bắc cầu.

Ba tính chất cần thiết của một quan hệ thứ tự đã được chứng minh, tức là, điều phải chứng minh.

\stepSubexercise Nếu $\left[\begin{aligned}
   A = \varnothing \\
   B = \varnothing
\end{aligned}\right.$ Khi đó, dựa vào kết quả của những bài tập trước, dễ thấy hàm số duy nhất từ một tập con nào đó của $A$ sang $B$ là $f_\varnothing = \varnothing$. Do đó, phần tử duy nhất trong $C$ là $(\varnothing; \varnothing)$.

Gọi tập hợp chứa các phần tử cực đại là $P_{\text{li}}$. Tập này có nghĩa do tiên đề tách lấy từ $C$. Xét $(A; f_A)$ ($f_A$ là một ánh xạ từ $A$ sang $B$) và gọi nó là $k$. $k \in P_{\text{li}}$ là do giả thiết cách đặt $P_{\text{li}}$. Nếu như có phần tử $l$ cũng thuộc $P_{\text{li}}$ sao cho $k\;R\;l$, tồn tại $j$ và $o$ sao cho 
\begin{itemize}
   \item $j \in \mathfrak{P}(A)$,
   \item $o$ là một hàm số từ $A$ sang $B$,
   \item $l = (j; o)$.
\end{itemize}

Bởi $j \in \mathfrak{P}(A)$, nhưng lại có $A \subseteq j$ do $k\;R\;l$ nên $j = A$. Khi đó, 
$$
   \forall u \in A, f_A(u) = o(u).
$$
$f_A$ và $o$ là hai ánh xạ có cùng tập xác định là $A$ và tập đích là $B$, cho nên $f_A = o$.

Vậy, $k = l$. Và từ đó, chúng ta có $k \in P_{\text{li}}$ theo định nghĩa phần tử cực đại.

Chúng ta cũng sẽ khẳng định rằng mọi phần tử không thể viết dưới dạng $(A; f_A)$ thì không phải là phần tử cực đại. Thật vậy, nếu như có $u \in C$ sao cho không tồn tại $g_A: A \to B$ để $u = (A; g_A)$, $u$ sẽ có dạng $(v; w)$ với $v \neq A$ (do nếu $v = A$ thì $w$ sẽ cần phải là một ánh xạ từ $A$ đến $B$ để có thể thuộc $C$). 
